\section{A uniform analytic price for the bare plaquette energy}
\label{sec:ex-wilson-vacuum-price}

The path construction uses uniform analytic bounds for the unperturbed Wilson
energy, before it compares that energy with a general path. This section
records the finite-volume estimate and its proof. All constants are independent
of the torus side.

Write $\mathcal A$ for the finite set of positive oriented plaquette types,
$N_{\mathcal A}=|\mathcal A|$, and let $d_G$ be the dimension of the fixed
matrix representation. For a link field $U$ on a torus of side $m$, put
\[
 W_m(U)=\sum_{x\in(\mathbb Z/m\mathbb Z)^4}\sum_{a\in\mathcal A}
       \left(1-\frac{\operatorname{Re}\operatorname{tr}U(\partial(x,a))}{d_G}\right).
\]
Here $U(\partial(x,a))$ is the ordered four-edge holonomy. The estimate below
is a coefficientwise price uniform in $m$; the action itself is extensive.
We use Lie-algebra coordinates $z$ at the identity and write
$\widetilde W_m(z)=W_m(\exp z)$; every derivative below means the derivative
of this coordinate function at $z=0$. The norm on link-coordinate arrays and
on their local Fourier-mode restrictions is the supremum norm.

For $t\in\mathbb N$, set $n=t+1$. A coefficient picture of order $t$ assigns a
function of $t$ independent complex momenta $q_1,\ldots,q_t$ to every list
of $n$ colors and oriented-link directions; the zeroth momentum is fixed by
$q_0=-\sum_{j=1}^tq_j$. The width-one strip is the set where every
coordinate of each independent momentum has imaginary part of absolute value
at most one. Let
$\operatorname{CTax}(R,0)=1$ and
$\operatorname{CTax}(R,k+1)=\operatorname{CTax}(R,k)/(R/2^{k+1})$.

\begin{proposition}[Uniform bare-energy coefficient price]
\label{prop:ex-wilson-stream}
There are constants $R_*>0$ and $A_*<\infty$, depending only on the compact
gauge representation and plaquette convention, and for each $t$ there is a
single picture $\mathcal W_t$ that works for every torus side $m$. Its
components are holomorphic on the width-one momentum strip and satisfy
\[
 \sup_{q\text{ in the strip}}|\mathcal W_t(q)|\leq
 A_t:=N_{\mathcal A}A_*\exp\bigl(16(t+1)^2\bigr)
       \frac{\operatorname{CTax}(R_*,t+1)}{(t+1)!}.
\]
Moreover, let $p_0,\ldots,p_t$ be dual-torus momenta, $c_0,\ldots,c_t$
colors, and $b_0,\ldots,b_t$ oriented-link directions. Let $e_{p,b,c}$
denote the corresponding Fourier-mode variation of the link field, and put
$q=(p_1,\ldots,p_t)$. Then
\[
 \frac{1}{n!}D^n \widetilde W_m(0)\bigl[e_{p_0,b_0,c_0},\ldots,e_{p_t,b_t,c_t}\bigr]
 =\begin{cases}
 |(\mathbb Z/m\mathbb Z)^4|\,\mathcal W_t(\mathbf c,\mathbf b;q),
       &\sum_{j=0}^t p_j=0,\\
 0,&\sum_{j=0}^t p_j\ne0.
 \end{cases}
\]
Thus the volume factor in the zero-total-momentum coefficient is exactly the
torus cardinality, and every nonzero-total-momentum coefficient vanishes.
\end{proposition}

\begin{proof}
First construct a single plaquette extension. Let $\mathfrak g_\mathbb C$ be
the complexification of the represented Lie algebra and let
$\mathcal E:\mathfrak g_\mathbb C\to\mathrm{Mat}_{d_G}(\mathbb C)$ be the
linear representation map. For a complex link-coordinate array $z$, the two
oriented link matrices are $\exp(\mathcal E z_e)$ and
$\exp(-\mathcal E z_e)$. These depend holomorphically on $z$, since the
matrix exponential is entire. Multiplication along the four edges, trace,
and the average of the forward and reverse loop traces are holomorphic too.
On real coordinates that average is the real trace of the holonomy, because
for the unitary holonomy $H$ one has
$\operatorname{tr}(H^{-1})=\overline{\operatorname{tr}(H)}$. Thus this
extension restricts to the real plaquette energy summand.

Continuity of the exponential at zero gives $R_*>0$ such that complex link
coordinates in the radius-$R_*$ ball give oriented matrices within
$\operatorname{rad}(1)$ of their physical unitary matrices. In the source's
Frobenius norm, this implies the correct product estimate as follows. If $M$
is within $r=\operatorname{rad}(1)$ of a unitary $V$, then for every matrix
$X$,
\[
 \|MX\|_F\leq\|VX\|_F+\|(M-V)X\|_F
       \leq(1+r)\|X\|_F,
\]
since left multiplication by $V$ is an isometry and the operator norm of
$M-V$ is at most its Frobenius norm. Iterating this inequality through four
factors, starting with $X=I$, gives
$\|M_1M_2M_3M_4\|_F\leq(1+r)^4\|I\|_F$. The reverse loop has the same
bound. The continuous pulse
\[
 (M,N)\longmapsto1-\frac{\operatorname{tr}M+\operatorname{tr}N}{2d_G}
\]
is bounded on the product of the two closed Frobenius balls of radius
$T=(1+r)^4\|I\|_F$ about zero, since both four-edge products lie there.
Compactness in finite dimension gives a common bound $A_*$.
Together these facts give a uniform holomorphic extension bounded by $A_*$
on a fixed ball, at every plaquette anchor and every torus size.

Each plaquette is supported in the radius-four coordinate box about its
anchor: every link in its four-edge boundary has an anchor at max-distance at
most four. Hence its real local action is unchanged when the cover
configuration is changed outside that box.

To define the common local symbol, place each radius-four box in the auxiliary
torus $\Lambda_9=(\mathbb Z/9\mathbb Z)^4$, with $9=2d+1$ and $d=4$. The
box embeds in this torus without identifications. For each plaquette type,
evaluate its local Taylor form there on the restricted Fourier modes
$y\mapsto {\bf1}_{(b,c)}\exp(i\langle p,\operatorname{lift}y\rangle)$,
and sum these local symbols over the plaquette types; this defines
$\mathcal W_t$ at the corresponding independent momenta. For a period of any
side $m$, pull a period-$m$ field $z$ to the auxiliary torus by
\[
 (P_mz)(y,b,c)=z(\operatorname{wrap}_m(\operatorname{lift}y),b,c).
\]
Then the local plaquette function satisfies
$\operatorname{local}_m(z)=\operatorname{local}_9(P_mz)$. The pullback of a
period-$m$ Fourier mode is exactly the restricted Fourier mode used above.
This identity holds even when $m<9$ and distinct auxiliary sites wrap to the
same period-$m$ link: evaluating $\operatorname{local}_9$ on $P_mz$ retains
all repeated-link contributions. Thus the local symbol is the same for every
period, including the small periods.
The picture variables are momenta, rather than link coordinates. To obtain
the strip bound, evaluate the order-$n$ Taylor multilinear form of the local
plaquette extension on the Fourier-mode vectors $e_{p_j,b_j,c_j}$. The Cauchy
derivative estimate on successively shrinking balls of radii $R_*/2,R_*/2^2,\ldots,R_*/2^n$ bounds its order-$n$
derivative by $A_*\operatorname{CTax}(R_*,n)$. The Taylor coefficient is
normalized by $1/n!$.
For a momentum in the width-one strip, a Fourier mode restricted to the
radius-four stencil has supremum norm at most $\exp(16n)$: the dependent momentum has imaginary coordinates
bounded by $t$, and a stencil displacement has four coordinates, each of
absolute value at most four. There are $n$ input modes in the multilinear form, so their
product costs at most $\exp(16n^2)$. This proves the claimed bound for each
plaquette type. Summing over the finite set $\mathcal A$ multiplies it by
$N_{\mathcal A}$ and yields a single picture $\mathcal W_t$, independent of
the torus side.

It remains to check the finite-volume coefficient identity. Expand the
translated plaquette sum and apply the multilinear coefficient to an ordered
momentum/color/direction tuple. Translating a local plaquette coefficient
by an anchor multiplies its Fourier mode by the character at total momentum
$\sum_{j=0}^t p_j$. The sum over all anchors is zero unless that total momentum is
zero; in that case each anchor contributes the same local coefficient, so
the sum is exactly $|(\mathbb Z/m\mathbb Z)^4|$ times the local symbol
$\mathcal W_t$ evaluated at the output phase. This proves the displayed
identity, including its volume normalization.
\end{proof}

\begin{remark}[Use in the selected path estimate]
\label{rem:ex-wilson-stream-use}
The source's all-orders summation lemma applies this same local-extension
argument to a countable family of interaction atoms, with the summability
condition
\[
 \sum_j b_j\exp\bigl((t+1)^2\sigma\,4d_j\bigr)<\infty,
\]
where $d_j$ is the support length and $b_j$ its absolute coefficient times
its exponential length weight. For the bare action the family is finite in
plaquette type, has support length four, and the preceding bound gives its
summability directly. Applying the lemma at orders three and four yields the
two Wilson reference pictures used in the one-step path estimate. They are
uniform over torus sides; no limit in the volume is used here.
\end{remark}
